\chapter{Introduction}\label{chap:introduction}

Term rewriting provides a small but expressive model of computation.  A term
rewrite system describes a computation by rules that replace matching subterms,
and a computation is obtained by applying such rules repeatedly.  Termination
asks whether this process can continue forever.  Besides being a fundamental
property in its own right, termination is a prerequisite for using rewrite
systems as executable specifications and for reasoning about the result of a
computation.  It has therefore been studied extensively, both as a mathematical
problem and as a target of automated verification; see, for example,
\cite{baader_nipkow_1999}.

Probabilistic term rewrite systems (PTRSs) extend this model by allowing a rule
to choose among several right-hand sides according to a probability
distribution.  They can represent randomised algorithms and probabilistic
program fragments while retaining the syntax and compositional structure of
term rewriting.  For such systems, classical termination is often unnecessarily
strong.  A probabilistic computation may admit an infinite run and nevertheless
terminate with probability one.  This motivates \emph{almost-sure termination}
(AST), which requires the probability of eventually reaching a normal form to
be one \cite{avanzini_dallago_yamada_2020}.  Importantly, AST is a statement
about probability mass rather than about every individual path.

\section{Motivation and Problem Statement}

The order in which redexes are contracted can change termination behaviour.  A
redex is an occurrence of a subterm that matches the left-hand side of a rule.
Under \emph{outermost} rewriting, a redex may be contracted only if it is not
strictly contained in another redex.  Under \emph{innermost} rewriting, the
chosen redex may not contain another redex.  These strategies correspond to two
opposite evaluation disciplines: outermost rewriting gives priority to the
surrounding computation, whereas innermost rewriting evaluates arguments first.

The difference is illustrated by the PTRS in
\cref{fig:probabilistic-strategy-gap}.  Its producer constructs a pair whose
first component is chosen probabilistically and whose second component invokes
the producer again.  Starting with a projection of the producer, outermost
evaluation discards the recursive component and returns the first component.
Innermost evaluation must instead keep expanding the recursive component before
it may apply the projection.  The same start term therefore terminates under
outermost rewriting but not under innermost rewriting.

\begin{figure}[htbp]
  \centering
  \begin{tikzpicture}[
    >={Latex[length=2mm]},
    rewritearrow/.style={->,semithick},
    term/.style={
      draw=black!45,
      rounded corners=1.5pt,
      fill=white,
      inner xsep=4pt,
      inner ysep=3pt,
      font=\scriptsize,
      align=center
    },
    rulebox/.style={
      draw=black!60,
      rounded corners=2pt,
      fill=black!2,
      inner xsep=6pt,
      inner ysep=4pt,
      font=\footnotesize,
      align=center
    },
    note/.style={font=\scriptsize,align=center}
  ]
    \node[rulebox,text width=12.1cm] at (0,4.35) {%
      $\begin{gathered}
        \P_{\mathsf{ex}}:\qquad
        \mathsf{produce}
        \to
        \left\{
          \tfrac12:\mathsf{pair}(\mathsf a,\mathsf{produce}),
          \tfrac12:\mathsf{pair}(\mathsf b,\mathsf{produce})
        \right\},
        \\
        \mathsf{first}(\mathsf{pair}(x,y))
        \to \{1:x\}.
      \end{gathered}$
    };

    \filldraw[
      rounded corners=3pt,
      fill=RoyalBlue!4,
      draw=RoyalBlue!55
    ] (-6.45,3.55) rectangle (-0.15,-1.65);
    \filldraw[
      rounded corners=3pt,
      fill=BurntOrange!4,
      draw=BurntOrange!65
    ] (0.15,3.55) rectangle (6.45,-1.65);

    \node[font=\small\bfseries\sffamily,text=RoyalBlue!75!black]
      at (-3.30,3.25) {Outermost};
    \node[font=\small\bfseries\sffamily,text=BurntOrange!85!black]
      at (3.30,3.25) {Innermost};

    \node[term] (o0) at (-3.30,2.75)
      {$\mathsf{first}(\mathsf{produce})$};
    \node[term,text width=5.05cm] (o1) at (-3.30,1.45)
      {$\mathsf{first}(\mathsf{pair}(c_1,\mathsf{produce}))$};
    \node[term] (o2) at (-3.30,0.20)
      {$c_1$ \textup{(normal form)}};

    \draw[rewritearrow] (o0) -- node[
      midway,anchor=east,xshift=-2pt,note
    ] {expand producer} (o1);
    \draw[rewritearrow] (o1) -- node[
      midway,anchor=west,xshift=2pt,note
    ] {contract root} (o2);
    \node[note,text=RoyalBlue!75!black] at (-3.30,-1.10) {
      the recursive argument is discarded\\
      termination probability: $1$
    };

    \node[term] (i0) at (3.30,2.75)
      {$\mathsf{first}(\mathsf{produce})$};
    \node[term,text width=5.05cm] (i1) at (3.30,1.65)
      {$\mathsf{first}(\mathsf{pair}(c_1,\mathsf{produce}))$};
    \node[term,text width=5.35cm] (i2) at (3.30,0.45)
      {$\mathsf{first}(\mathsf{pair}(c_1,
        \mathsf{pair}(c_2,\mathsf{produce})))$};
    \node[term,text width=5.35cm] (i3) at (3.30,-0.65)
      {$\mathsf{first}(\mathsf{pair}(c_1,\ldots,
        \mathsf{pair}(c_n,\mathsf{produce})\ldots))$};

    \draw[rewritearrow] (i0) -- node[
      midway,anchor=east,xshift=-2pt,note
    ] {expand producer} (i1);
    \draw[rewritearrow] (i1) -- node[
      midway,anchor=west,xshift=2pt,note
    ] {expand tail} (i2);
    \draw[rewritearrow,dashed] (i2) -- node[
      midway,anchor=west,xshift=2pt,note
    ] {repeat} (i3);
    \node[note,text=BurntOrange!85!black] at (3.30,-1.35) {
      the root projection remains blocked\\
      termination probability: $0$
    };
  \end{tikzpicture}
  \caption{Strategy-dependent behaviour from
    $\mathsf{first}(\mathsf{produce})$.  Each $c_i$ is independently chosen as
    $\mathsf a$ or $\mathsf b$, each with probability $\tfrac12$.}
  \label{fig:probabilistic-strategy-gap}
\end{figure}

A formal account using lifted probabilistic derivations is given in
\cref{ex:probabilistic-strategy-difference}.

\Needspace{10\baselineskip}
Automated proof techniques are especially well developed for innermost
rewriting.  This suggests reducing an outermost termination problem to an
innermost one.  The central question of this thesis is therefore:

\begin{quote}
Can outermost almost-sure termination of a PTRS be established by transforming
it into a PTRS whose innermost almost-sure termination can be analysed instead?
\end{quote}

The intended transformation must do more than reproduce the set of reachable
terms.  Soundness requires it to preserve the probabilities of the original
alternatives and to represent nondeterministic choices between redexes and
rules.  Completeness additionally requires control over computations beginning
with arbitrary terms over the extended signature, including administrative
configurations.  Such computations must not introduce positive non-termination
probability that is absent from the source system.  The essential new issue in
the probabilistic setting is therefore the preservation of probability mass,
rather than merely the existence of infinite reductions.

\section{Related Work}

Transformation-based proofs of strategy-sensitive termination have a history in
ordinary term rewriting.  Raffelsieper and Zantema transform an outermost
termination problem into an unrestricted termination problem.  Their principal
transformation is sound for outermost ground termination and is effectively
constructible as a finite TRS for finite quasi-left-linear systems
\cite{raffelsieper_zantema_2009}.  Thiemann later introduced a different marker
transformation for finite TRSs and proved the exact correspondence
\[
  \R \text{ is outermost terminating}
  \quad\Longleftrightarrow\quad
  \R' \text{ is innermost terminating}.
\]
This construction is the deterministic starting point for the transformation
studied in this thesis \cite{thiemann_2009_outermost}.

Probabilistic rewriting requires a separate semantic treatment.  Avanzini,
Dal~Lago, and Yamada give a systematic account of PTRSs and probabilistic
termination notions \cite{avanzini_dallago_yamada_2020}.  Kassing and Giesl
study when innermost results for a PTRS can be transferred to unrestricted
probabilistic rewriting, and show that additional syntactic restrictions are
needed because duplication can affect probability mass
\cite{kassing_giesl_2025}.  Their work concerns innermost versus full rewriting;
it does not provide the outermost-to-innermost transformation considered here.

The deterministic transformation suggests the construction, but its
correctness proof does not transfer automatically to PTRSs.  In a deterministic
system, one infinite reduction is enough to refute termination.  In a
probabilistic system, the existence of infinite paths is compatible with AST
when the total probability of infinite computations is zero.  Consequently,
correctness has to be proved for infinite lifted derivations, or equivalently
for fully evaluated rewrite sequence trees, while retaining their branch
probabilities.

\section{Contribution}

This thesis adapts the marker-based outermost-to-innermost transformation of
\cite{thiemann_2009_outermost} to probabilistic term rewrite systems.  Given a
PTRS $\P$, the construction produces a PTRS $\P'$ over an extended signature.
Fresh administrative symbols guide an innermost computation to positions that
represent permitted outermost steps of $\P$.  The rules that move these symbols
are deterministic; probabilistic branching occurs only when an original rule
is simulated.

The main contributions are:

\begin{enumerate}
  \item a probabilistic outermost-to-innermost transformation, including the
        extended signature and its administrative rewrite rules;
  \item a simulation result showing that every outermost probabilistic step is
        represented by a finite, non-empty sequence of innermost steps while
        preserving the probabilities of all outcomes; and
  \item a lifting of the simulation and projection arguments to infinite
        probabilistic derivations, yielding the sound-and-complete equivalence
        \[
          \AST[\oto_{\P}]
          \quad\Longleftrightarrow\quad
          \AST[\ito_{\P'}].
        \]
\end{enumerate}

Both sides quantify over all terms, including non-ground terms.  In particular,
the target property concerns every term over the extended signature, not only
terms produced by the designated source encoding.  The result established here
concerns AST; no corresponding claim for PAST or SAST is made.

\section{Structure of the Thesis}

The remainder of this thesis is organised as follows.
\Cref{chap:preliminaries} introduces terms, rewrite strategies, finite
multi-distributions, probabilistic rewriting, rewrite sequence trees, and the
termination notions used throughout the thesis.  It also recalls the
deterministic transformation result on which the new construction is based.
\Cref{chap:main-chapter} presents the probabilistic transformation and proves
its simulation and correctness properties.  \Cref{chap:evaluation} illustrates
the construction on representative examples and discusses its applicability
and limitations.  Finally, \cref{chap:conclusion} summarises the results and
outlines possible extensions.
